\documentclass[10pt,a4paper]{amsart}
\usepackage{amsmath,amssymb,amsthm,graphicx}
\usepackage{verbatim}
\newenvironment{reference}{\comment}{\endcomment}
%\newenvironment{reference}{}{}
\newenvironment{slogan}{\comment}{\endcomment}
\newenvironment{history}{\comment}{\endcomment}

% For commutative diagrams we use Xy-pic
\usepackage[all]{xy}

% We use 2cell for 2-commutative diagrams.
\xyoption{2cell}
% Theorem environments.
%
\theoremstyle{plain}
\newtheorem{theorem}[subsection]{Theorem}
\newtheorem{proposition}[subsection]{Proposition}
\newtheorem{lemma}[subsection]{Lemma}

\theoremstyle{definition}
\newtheorem{definition}[subsection]{Definition}
\newtheorem{example}[subsection]{Example}
\newtheorem{exercise}[subsection]{Exercise}
\newtheorem{situation}[subsection]{Situation}

\theoremstyle{remark}
\newtheorem{remark}[subsection]{Remark}
\newtheorem{remarks}[subsection]{Remarks}

\numberwithin{equation}{subsection}

% Macros
%
\def\lim{\mathop{\mathrm{lim}}\nolimits}
\def\colim{\mathop{\mathrm{colim}}\nolimits}
\def\Spec{\mathop{\mathrm{Spec}}}
\def\Hom{\mathop{\mathrm{Hom}}\nolimits}
\def\Ext{\mathop{\mathrm{Ext}}\nolimits}
\def\SheafHom{\mathop{\mathcal{H}\!\mathit{om}}\nolimits}
\def\SheafExt{\mathop{\mathcal{E}\!\mathit{xt}}\nolimits}
\def\Sch{\mathit{Sch}}
\def\Mor{\mathop{\mathrm{Mor}}\nolimits}
\def\Ob{\mathop{\mathrm{Ob}}\nolimits}
\def\Sh{\mathop{\mathit{Sh}}\nolimits}
\def\NL{\mathop{N\!L}\nolimits}
\def\CH{\mathop{\mathrm{CH}}\nolimits}
\def\proetale{{pro\text{-}\acute{e}tale}}
\def\etale{{\acute{e}tale}}
\def\QCoh{\mathit{QCoh}}
\def\Ker{\mathop{\mathrm{Ker}}}
\def\Im{\mathop{\mathrm{Im}}}
\def\Coker{\mathop{\mathrm{Coker}}}
\def\Coim{\mathop{\mathrm{Coim}}}

% Boxtimes
%
\DeclareMathSymbol{\boxtimes}{\mathbin}{AMSa}{"02}

%
% Macros for moduli stacks/spaces
%
\def\QCohstack{\mathcal{QC}\!\mathit{oh}}
\def\Cohstack{\mathcal{C}\!\mathit{oh}}
\def\Spacesstack{\mathcal{S}\!\mathit{paces}}
\def\Quotfunctor{\mathrm{Quot}}
\def\Hilbfunctor{\mathrm{Hilb}}
\def\Curvesstack{\mathcal{C}\!\mathit{urves}}
\def\Polarizedstack{\mathcal{P}\!\mathit{olarized}}
\def\Complexesstack{\mathcal{C}\!\mathit{omplexes}}
% \Pic is the operator that assigns to X its picard group, usage \Pic(X)
% \Picardstack_{X/B} denotes the Picard stack of X over B
% \Picardfunctor_{X/B} denotes the Picard functor of X over B
\def\Pic{\mathop{\mathrm{Pic}}\nolimits}
\def\Picardstack{\mathcal{P}\!\mathit{ic}}
\def\Picardfunctor{\mathrm{Pic}}
\def\Deformationcategory{\mathcal{D}\!\mathit{ef}}
\usepackage{cleveref}
\crefname{originallemma}{original lemma}{original lemmas}
\Crefname{originallemma}{Original lemma}{Original lemmas}
\makeatletter
\newlabel{algebra-lemma-completion-finite-extension}{{07N9}{source}{}{sourcepointer-1}{}}
\newlabel{algebra-lemma-completion-finite-extension@cref}{{[originallemma][0][]07N9}{[1][0][]source}{}{}{}}
\makeatother
\crefname{originallemma}{original lemma}{original lemmas}
\Crefname{originallemma}{Original lemma}{Original lemmas}
\makeatletter
\newlabel{algebra-lemma-descent-regular}{{07NG}{source}{}{sourcepointer-2}{}}
\newlabel{algebra-lemma-descent-regular@cref}{{[originallemma][0][]07NG}{[1][0][]source}{}{}{}}
\makeatother
\crefname{originallemma}{original lemma}{original lemmas}
\Crefname{originallemma}{Original lemma}{Original lemmas}
\makeatletter
\newlabel{coherent-lemma-flat-base-change-cohomology}{{02KH}{source}{}{sourcepointer-3}{}}
\newlabel{coherent-lemma-flat-base-change-cohomology@cref}{{[originallemma][0][]02KH}{[1][0][]source}{}{}{}}
\makeatother
\crefname{originallemma}{original lemma}{original lemmas}
\Crefname{originallemma}{Original lemma}{Original lemmas}
\makeatletter
\newlabel{coherent-lemma-proper-over-affine-cohomology-finite}{{02O6}{source}{}{sourcepointer-4}{}}
\newlabel{coherent-lemma-proper-over-affine-cohomology-finite@cref}{{[originallemma][0][]02O6}{[1][0][]source}{}{}{}}
\makeatother
\crefname{originalreference}{original reference}{original references}
\Crefname{originalreference}{Original reference}{Original references}
\makeatletter
\newlabel{lemma-algebra-helper}{{0BGH}{source}{}{sourcepointer-5}{}}
\newlabel{lemma-algebra-helper@cref}{{[originalreference][0][]0BGH}{[1][0][]source}{}{}{}}
\makeatother
\crefname{originalreference}{original reference}{original references}
\Crefname{originalreference}{Original reference}{Original references}
\makeatletter
\newlabel{lemma-port-regularity-to-completion}{{0BG6}{source}{}{sourcepointer-6}{}}
\newlabel{lemma-port-regularity-to-completion@cref}{{[originalreference][0][]0BG6}{[1][0][]source}{}{}{}}
\makeatother
\crefname{originalreference}{original reference}{original references}
\Crefname{originalreference}{Original reference}{Original references}
\makeatletter
\newlabel{lemma-regular-alteration-implies}{{0BGG}{source}{}{sourcepointer-7}{}}
\newlabel{lemma-regular-alteration-implies@cref}{{[originalreference][0][]0BGG}{[1][0][]source}{}{}{}}
\makeatother
\crefname{originallemma}{original lemma}{original lemmas}
\Crefname{originallemma}{Original lemma}{Original lemmas}
\makeatletter
\newlabel{properties-lemma-normal-integral-sections}{{0358}{source}{}{sourcepointer-8}{}}
\newlabel{properties-lemma-normal-integral-sections@cref}{{[originallemma][0][]0358}{[1][0][]source}{}{}{}}
\makeatother
\crefname{originallemma}{original lemma}{original lemmas}
\Crefname{originallemma}{Original lemma}{Original lemmas}
\makeatletter
\newlabel{properties-lemma-regular-normal}{{0569}{source}{}{sourcepointer-9}{}}
\newlabel{properties-lemma-regular-normal@cref}{{[originallemma][0][]0569}{[1][0][]source}{}{}{}}
\makeatother
\title{Consequences of the existence of a regular alteration for a local Noetherian domain}
\author{The Stacks Project Authors}
\date{stacks-project:resolve.tex:local-snapshot-20261012; GNU FDL 1.2 or later}
\begin{document}
\maketitle
\noindent\textit{Editorial note.} All delivered mathematical text is the author's own, copied from the pinned Stacks Project chapter file resolve.tex; no statement or proof step was rewritten, reordered or removed.\par\medskip
\noindent\textit{Editorial note.} The counted proof is the author's complete argument: the use of the ring of global sections of the alteration to obtain finiteness of the integral closure and to reduce the normality assertion to a completion, the reduction to the case where the global sections equal the base ring with the connectedness and normality argument giving a normal domain, and the deduction of regularity of the generic formal fibre by descent along a finite faithfully flat map.\par\medskip
\noindent\textit{Editorial note.} The label lemma-regular-alteration-implies-local names the counted statement, and the proof environment immediately following it in the source is delivered verbatim as its proof.\par\medskip
\noindent\textit{Editorial note.} The counted statement is carried by the author's own lemma environment in the chapter source.\par\medskip
\noindent\textit{Editorial note.} Pre-open gate: python operations/agent-authoring/tools/precard\_gate.py --source resolve.tex --label lemma-regular-alteration-implies-local --start 3826 --end 3898 --local-id 28075 returned RESULT=OK (L28075 free, occupancy FREE, public interval clear). Pre-assembly tag self-check: canonical tag resolved and not present in reference/public-items.jsonl.\par\medskip
\noindent\textit{Editorial note.} The statement and proof are copied from the pinned chapter file; the cited inputs are kept as external pointers with their canonical Stacks tags.\par\medskip
\noindent\textit{Allowed standard prerequisites (editorial).}\begin{itemize}
\item Existence of a nonempty open and finiteness of the integral closure for an alteration. Same chapter, the lemma deducing the existence of an open on which the ring is regular and the finiteness of the integral closure from the existence of an alteration. Supplies the first two assertions of the lemma and sets up the notation reused in the proof.
\item Finiteness of the global sections of a proper scheme over an affine base. Cohomology of Schemes, the lemma that the global sections of a proper scheme over an affine base form a finite module, together with the flat base change lemma used twice. Supplies the finiteness of the ring of global sections and the identification of the global sections after flat base change.
\item Normality of the global sections of a normal scheme. Properties, the lemmas that a regular scheme is normal and that the global sections of a normal integral scheme form a normal domain. Supplies the normality of the ring of global sections and the conclusion that the completed ring is a normal domain.
\item Regularity of the source after base change to the completion, and descent of regularity. Same chapter, the lemma transferring regularity of an alteration along the completion, and Algebra, the lemma descending regularity along a finite faithfully flat map. Supply the regularity of the base changed alteration and the conclusion about the generic formal fibre.
\end{itemize}\medskip
\ifdefined\hypertarget\hypertarget{sourcepointer-1}{}\fi
\noindent\textit{Source reference.} Original lemma 07N9: sources/stacks-chow/tags.txt, source lines 9754--9754. Prerequisite pointer invoked at proof line(s) 3856; it names a standard fact of a sibling chapter that the counted argument uses, and every step of the counted proof built on it is retained verbatim in the delivered body.\par\medskip
\ifdefined\hypertarget\hypertarget{sourcepointer-2}{}\fi
\noindent\textit{Source reference.} Original lemma 07NG: sources/stacks-chow/tags.txt, source lines 9761--9761. Prerequisite pointer invoked at proof line(s) 3897; it names a standard fact of a sibling chapter that the counted argument uses, and every step of the counted proof built on it is retained verbatim in the delivered body.\par\medskip
\ifdefined\hypertarget\hypertarget{sourcepointer-3}{}\fi
\noindent\textit{Source reference.} Original lemma 02KH: sources/stacks-chow/tags.txt, source lines 3372--3372. Prerequisite pointer invoked at proof line(s) 3866; it names a standard fact of a sibling chapter that the counted argument uses, and every step of the counted proof built on it is retained verbatim in the delivered body.\par\medskip
\ifdefined\hypertarget\hypertarget{sourcepointer-4}{}\fi
\noindent\textit{Source reference.} Original lemma 02O6: sources/stacks-chow/tags.txt, source lines 3505--3505. Prerequisite pointer invoked at proof line(s) 3845; it names a standard fact of a sibling chapter that the counted argument uses, and every step of the counted proof built on it is retained verbatim in the delivered body.\par\medskip
\ifdefined\hypertarget\hypertarget{sourcepointer-5}{}\fi
\noindent\textit{Source reference.} Original reference 0BGH: sources/stacks-resolve/pinned-original-source.tex, source lines 3776--3776. Prerequisite pointer invoked at proof line(s) 3851; it names a standard fact of the same chapter that the counted argument uses, and every step of the counted proof built on it is retained verbatim in the delivered body.\par\medskip
\ifdefined\hypertarget\hypertarget{sourcepointer-6}{}\fi
\noindent\textit{Source reference.} Original reference 0BG6: sources/stacks-resolve/pinned-original-source.tex, source lines 2997--2997. Prerequisite pointer invoked at proof line(s) 3873; it names a standard fact of the same chapter that the counted argument uses, and every step of the counted proof built on it is retained verbatim in the delivered body.\par\medskip
\ifdefined\hypertarget\hypertarget{sourcepointer-7}{}\fi
\noindent\textit{Source reference.} Original reference 0BGG: sources/stacks-resolve/pinned-original-source.tex, source lines 3747--3747. Prerequisite pointer invoked at proof line(s) 3841; it names a standard fact of the same chapter that the counted argument uses, and every step of the counted proof built on it is retained verbatim in the delivered body.\par\medskip
\ifdefined\hypertarget\hypertarget{sourcepointer-8}{}\fi
\noindent\textit{Source reference.} Original lemma 0358: sources/stacks-chow/tags.txt, source lines 4119--4119. Prerequisite pointer invoked at proof line(s) 3849; it names a standard fact of a sibling chapter that the counted argument uses, and every step of the counted proof built on it is retained verbatim in the delivered body.\par\medskip
\ifdefined\hypertarget\hypertarget{sourcepointer-9}{}\fi
\noindent\textit{Source reference.} Original lemma 0569: sources/stacks-chow/tags.txt, source lines 6702--6702. Prerequisite pointer invoked at proof line(s) 3847; it names a standard fact of a sibling chapter that the counted argument uses, and every step of the counted proof built on it is retained verbatim in the delivered body.\par\medskip
\begin{lemma}
\label{lemma-regular-alteration-implies-local}
Let $(A, \mathfrak m, \kappa)$ be a local Noetherian domain.
Assume there exists an alteration $f : X \to \Spec(A)$
with $X$ regular. Then
\begin{enumerate}
\item there exists a nonzero $f \in A$ such that $A_f$ is regular,
\item the integral closure $B$ of $A$ in its fraction field is finite over $A$,
\item the $\mathfrak m$-adic completion of $B$ is a normal ring, i.e., the
completions of $B$ at its maximal ideals are normal domains, and
\item the generic formal fibre of $A$ is regular.
\end{enumerate}
\end{lemma}

\begin{proof}
Parts (1) and (2) follow from Lemma \ref{lemma-regular-alteration-implies}.
We have to redo part of the proof of that lemma in order to set up notation
for the proof of (3). Set $C = \Gamma(X, \mathcal{O}_X)$. By
Cohomology of Schemes, Lemma
\ref{coherent-lemma-proper-over-affine-cohomology-finite}
we see that $C$ is a finite $A$-module. As $X$ is normal
(Properties, Lemma \ref{properties-lemma-regular-normal})
we see that $C$ is normal domain
(Properties, Lemma \ref{properties-lemma-normal-integral-sections}).
Thus $B \subset C$ and we conclude that $B$ is finite over $A$
as $A$ is Noetherian. By Lemma \ref{lemma-algebra-helper}
in order to prove (3) it suffices to show
that the $\mathfrak m$-adic completion $C^\wedge$ is normal.

\medskip\noindent
By Algebra, Lemma \ref{algebra-lemma-completion-finite-extension}
the completion $C^\wedge$ is the product of the completions of
$C$ at the prime ideals of $C$ lying over $\mathfrak m$.
There are finitely many of these and these are the maximal
ideals $\mathfrak m_1, \ldots, \mathfrak m_r$ of $C$.
(The corresponding result for $B$ explains the final statement of the lemma.)
Thus replacing $A$ by $C_{\mathfrak m_i}$ and $X$ by
$X_i = X \times_{\Spec(C)} \Spec(C_{\mathfrak m_i})$
we reduce to the case discussed in the next paragraph.
(Note that $\Gamma(X_i, \mathcal{O}) = C_{\mathfrak m_i}$ by
Cohomology of Schemes, Lemma \ref{coherent-lemma-flat-base-change-cohomology}.)

\medskip\noindent
Here $A$ is a Noetherian local normal domain and $f : X \to \Spec(A)$ is a
regular alteration with $\Gamma(X, \mathcal{O}_X) = A$.
We have to show that the completion $A^\wedge$
of $A$ is a normal domain. By
Lemma \ref{lemma-port-regularity-to-completion}
$Y = X \times_{\Spec(A)} \Spec(A^\wedge)$ is regular.
Since $\Gamma(Y, \mathcal{O}_Y) = A^\wedge$ by
Cohomology of Schemes, Lemma \ref{coherent-lemma-flat-base-change-cohomology},
we conclude that $A^\wedge$ is normal as before.
Namely, $Y$ is normal by
Properties, Lemma \ref{properties-lemma-regular-normal}.
It is connected because $\Gamma(Y, \mathcal{O}_Y) = A^\wedge$ is local.
Hence $Y$ is normal and integral (as connected and normal
implies integral for Noetherian schemes). Thus
$\Gamma(Y, \mathcal{O}_Y) = A^\wedge$
is a normal domain by
Properties, Lemma \ref{properties-lemma-normal-integral-sections}.
This proves (3).

\medskip\noindent
Proof of (4). Let $\eta \in \Spec(A)$ denote the generic point
and denote by a subscript $\eta$ the base change to $\eta$.
Since $f$ is an alteration, the scheme $X_\eta$ is finite and
faithfully flat over $\eta$. Since $Y = X \times_{\Spec(A)} \Spec(A^\wedge)$
is regular by Lemma \ref{lemma-port-regularity-to-completion}
we see that $Y_\eta$ is regular (as a limit of opens in $Y$).
Then $Y_\eta \to \Spec(A^\wedge \otimes_A \kappa(\eta))$ is finite
faithfully flat onto the generic formal fibre. We conclude by
Algebra, Lemma \ref{algebra-lemma-descent-regular}.
\end{proof}

\end{document}
